\chapter{The design of a decoupled scientific engine}\label{ch:arch}

\begin{tcolorbox}[colback=gray!5,colframe=gray!50,title={In this chapter}]
Suppose a colleague asks for a new element, a simple bar, and wants to run it in the static solver, the modal solver and the nonlinear drivers. In a code where each solver knows which elements it handles, that means opening three solvers and adding a branch to each; the next element needs the same again. The two packages in this book, \fea{} and \rom{}, are built so that this is not necessary: a new element, material, geometry, load, solution strategy or linear-algebra backend is added by a subclass or a registry entry. This chapter shows the structure that allows this and tests the claims it rests on. It follows the flow from material to post-processing, the small interface through which solvers talk to elements, a new element of about twenty lines added to a solver that was not changed, an interface test of all 34 registered elements, the rule that \rom{} never imports \fea{}, and the backend dispatcher, whose cost on sparse systems we measured and the package documentation does not mention.\par\smallskip
After the chapter you should be able to add an element by subclassing and know which interface checks it has to pass, say why \rom{} does not import \fea{} and how that is checked, and explain what the \code{auto} backend setting costs on sparse systems. The signatures of the classes named here are in the \ghfile{API_REFERENCE.md}{API reference} of the repository, and the \ghfile{USER_GUIDE.md}{user guide} covers installation and first runs.
\end{tcolorbox}

\section{One system, many problems}\label{sec:ar-system}

Whatever the physical problem is (a membrane, a beam, a shell with contact), the solvers of \fea{} work on one algebraic form,
\begin{equation}\label{eq:ar-system}
\bm M\ddot{\bm u}+\bm C\dot{\bm u}+\bm K\bm u=\bm F ,
\end{equation}
or, for nonlinear problems, on the residual $\bm f_{int}(\bm u)-\bm F$ and its tangent (Chapter~\ref{ch:nonlin}). The vectors and matrices are plain NumPy or SciPy arrays. Everything that is specific to a problem is pushed into the construction of $\bm M,\bm C,\bm K$ (or of $\bm f_{int}$ and $\bm K_T$), and everything that is specific to a solution method is in the solvers. The question for this chapter is where the boundary between the two sits, and how narrow it can be made.

\begin{sidenote}{The matrices of the semi-discrete equation}
$\bm u$ collects the unknowns of the model, the displacements (and rotations) of all nodes, and dots denote time derivatives. $\bm M$ is the mass matrix, which turns accelerations into inertial forces, $\bm C$ the damping matrix, which turns velocities into dissipative forces, and $\bm K$ the stiffness matrix, which turns displacements into elastic forces; $\bm F$ is the vector of external loads. A static problem drops the first two terms and solves $\bm K\bm u=\bm F$. The matrix $\bm K$ is symmetric, \emph{sparse} (a node is coupled only to the nodes that share an element with it, so most entries are zero), and positive definite only once the supports have removed the rigid-body motions; for a free body it is singular, because the body can move without storing energy.
\end{sidenote}

\section{The flow and the layers}\label{sec:ar-flow}

\Cref{fig:ar-flow} shows the flow of a typical analysis. A \emph{material} gives the constitutive matrix (for example \code{D_plane_stress(Material(...))}); a \emph{mesh} gives nodes and element connectivity (\code{Mesh}, or \code{MultiBlockMesh} for several element topologies, Chapter~\ref{ch:mesh}); an \emph{element} formulation turns coordinates and a material into element matrices or forces; the system \anchor{fea_engine.solver.FESystem} assembles them, applies boundary conditions and loads, and calls a \emph{solver}; \emph{post-processing} recovers stress, strain and fields from the result. A layer knows only the layer next to it: the mesh does not know which element will use it, and the solver does not know which element it is solving for.

\tikzfigure[0.98\linewidth]{ch01_flow}{The flow of an analysis in \fea{} (blue: package code; grey: data; violet: \rom{}; green: validation only; dashed: optional). The link to \rom{} is plain arrays, never an import.}{fig:ar-flow}

\section{The element interface}\label{sec:ar-interface}

\begin{sidenote}{Registry and base class}
The \emph{registry} \code{ELEMENT\_REGISTRY} is a dictionary from a short name to an element class, so a mesh or an input file can name an element by string. The \emph{base class} \code{Element} supplies everything an element has in common, such as the Gauss loop, so a new element only states what is specific to it.
\end{sidenote}

The narrow boundary is the base class \anchor{fea_engine.elements.base.Element}. An element declares four class attributes, \code{n_nodes}, \code{dofs_per_node}, \code{dim} (the dimension of the natural-coordinate space) and \code{gauss_order}, and provides the shape functions in the natural coordinates, \code{shape_and_derivs}. From these the base class implements the Gauss loop \emph{once}: if the element supplies a strain-displacement matrix \code{B_matrix}, the generic \code{stiffness} computes $\int\bm B^{T}\bm D\bm B$ (\Cref{eq:el-K}) and the generic \code{mass} computes $\int\rho\bm N^{T}\bm N$. An element that needs more overrides the method it needs: the Mindlin plate overrides \code{stiffness} for selective reduced integration, and the Euler--Bernoulli beam uses a closed form. Nonlinear elements provide the pair \code{internal_force} and \code{tangent_stiffness} instead (or as well). Optional capabilities are declared by attributes and methods the solvers look for: \code{supports_batch} with the batched kernels of Chapter~\ref{ch:batch}, \code{state} and \code{commit_state} for history variables (Chapter~\ref{ch:nonlin}), and \code{lumped_mass} for the explicit method (Chapter~\ref{ch:dyn}).

The contact elements of \fea{} are in the same registry but have no interpolation: they define force and tangent from a gap function and do not implement \code{shape_and_derivs}, which raises \code{NotImplementedError} for them. The interface is therefore not a single fixed set of methods, but ``what the solver calls'': the assembly asks the element for matrices (linear) or for force and tangent (nonlinear), and the element is free in how it computes them.

\tikzfigure[0.62\linewidth]{ch01_interface}{The element interface. The base class owns the Gauss loop; a subclass supplies the interpolation and overrides only what differs.}{fig:ar-interface}

\begin{algbox}{Global assembly of $\bm K$ (and, identically, $\bm M$)}
\begin{pseudo}
\Require system with element blocks, material (or per-block materials) $\bm D$, node coordinates, element-to-dof maps
\State allocate $\bm K$ (dense, or sparse triplets)
\For{each block $b$ and each element $e$ of $b$}
  \State gather the node coordinates $\bm x_e$ and the dof map $\mathcal D_e$
  \State $\bm K_e\gets$ \texttt{element.stiffness}$(\bm x_e,\bm D)$ \Comment{Gauss loop of the base class, or an override}
  \State $\bm K[\mathcal D_e,\mathcal D_e]\gets\bm K[\mathcal D_e,\mathcal D_e]+\bm K_e$
\EndFor
\State apply boundary conditions on the free/fixed split afterwards
\Ensure assembled stiffness (calling again \emph{replaces} it unless \texttt{accumulate=True})
\end{pseudo}
\end{algbox}

Assembly costs one element matrix of $O(n_g\,d^3)$ per element plus a scatter-add, and for small elements the Python loop over the elements dominates; Chapter~\ref{ch:batch} removes it. Two things go wrong: an element with a non-positive Jacobian raises no error (Chapter~\ref{ch:mesh}), and an element that does not implement the requested method fails at the call, as a contact element does when asked for \code{shape_and_derivs}. The code is \code{assemble_stiffness} and \code{assemble_mass} of \anchor{fea_engine.solver.FESystem}, with the Gauss loop in \anchor{fea_engine.elements.base.Element}; \Cref{tab:ar-registry} and \Cref{lst:ar-bar} are the checks.

\section{Extending the engine: a new element}\label{sec:ar-extend}

A new element should not need any existing code to change. \Cref{lst:ar-bar} defines a two-node bar with one displacement per node (a one-dimensional element, $\bm B=\mathrm dN/\mathrm dx$, with the stiffness $\bm D=EA$ as a $1\times1$ matrix) as a subclass, registers it with one line (after \texttt{from fea\_engine.elements import Element, ELEMENT\_REGISTRY, jacobian}), and uses it with the existing assembly, static solver and modal solver.

\begin{lstlisting}[caption={A new element by subclassing, the one registry line, and a check that the package was not edited (run with the package on the path).},label={lst:ar-bar}]
import hashlib, pathlib, ast
import numpy as np
import fea_engine
from fea_engine import FESystem
from fea_engine.elements import Element, ELEMENT_REGISTRY, jacobian
from fea_engine.mesh import line_mesh

pkg = pathlib.Path(fea_engine.__file__).parent
digest = lambda: hashlib.sha256(b"".join(p.read_bytes() for p in sorted(pkg.rglob("*.py")))).hexdigest()
before = digest()

# what the solvers know about concrete elements: names of element classes used in their code (not comments)
names = {c.__name__ for c in ELEMENT_REGISTRY.values()} - {"Element"}
used = 0
for f in ["solver.py", "nonlinear_solver.py", "contact_solvers.py"]:
    for node in ast.walk(ast.parse((pkg / f).read_text())):
        used += isinstance(node, ast.Name) and node.id in names
        used += isinstance(node, ast.Attribute) and node.attr in names
print(f"concrete element classes named in the code of the three solver modules: {used}")

class Bar1D(Element):                                    # the new element
    n_nodes, dofs_per_node, dim, gauss_order = 2, 1, 1, 1
    dof_names = ("u",)

    def shape_and_derivs(self, natural_coords):
        (xi,) = natural_coords
        return np.array([(1 - xi) / 2, (1 + xi) / 2]), np.array([[-0.5, 0.5]])

    def B_matrix(self, natural_coords, elem_coords):
        _, dN = self.shape_and_derivs(natural_coords)
        J, detJ = jacobian(dN, elem_coords)
        return np.linalg.solve(J, dN), detJ

ELEMENT_REGISTRY["bar1d"] = Bar1D                        # the one registry line

E, A, rho, L, P = 210e9, 1.0e-4, 7800.0, 2.0, 1000.0
for n in (1, 4, 32):
    s = FESystem(line_mesh(L, n), ELEMENT_REGISTRY["bar1d"]())
    s.assemble_stiffness(np.array([[E * A]])); s.assemble_mass(np.array([[rho * A]]))
    s.fix_dofs([0], [0]); s.add_nodal_force([n], 0, P)
    U = s.solve_static(); f, _ = s.solve_modal(n_modes=1)
    f_exact = np.sqrt(E / rho) / (4 * L)
    print(f"{n:2d} elements: tip error {abs(U[n] / (P * L / (E * A)) - 1):.0e}, "
          f"first frequency error {abs(f[0] / f_exact - 1):.1e}")

print("package source changed:", digest() != before)
\end{lstlisting}

\begin{lstlisting}[numbers=none,language={},keywordstyle={},commentstyle={},frame=single,aboveskip=2pt]
concrete element classes named in the code of the three solver modules: 0
 1 elements: tip error 0e+00, first frequency error 2.7e-01
 4 elements: tip error 1e-15, first frequency error 1.3e-02
32 elements: tip error 0e+00, first frequency error 2.0e-04
package source changed: False
\end{lstlisting}

The naive design is the one we would have to write if solvers knew their elements, with a branch per element in each solver; here the three solver modules name no concrete element class, so there is nothing to edit, and the check of the package source confirms that the run changed nothing. The new element runs in the static and the modal solver, with a tip displacement that is exact (the linear element is exact for this problem) and a frequency error that falls as the mesh is refined.

The longer study behind the listing was run for 1, 2, 4, 8, 16 and 32 elements ($E=210$ GPa, $A=10^{-4}$ m$^2$, $\rho=7800$ kg/m$^3$, $L=2$ m). The static tip displacement equals $PL/(EA)$ to \cdeExtStatic{} (the linear element is exact for this problem), and the first axial frequency converges to $\sqrt{E/\rho}/(4L)$ with observed order \cdeExtFreqOrder{}, from an error of \cdeExtFreqCoarse{} with one element. The existing solvers ran on the new element without any change. That does not mean every solver works with every new element: a path-dependent or batched element also has to meet the conventions of \Cref{sec:ar-interface}.

\section{The six extension points}\label{sec:ar-points}

\Cref{tab:ar-points} lists where each kind of extension enters. In each case the repository's own convention is a dictionary or a base class, so that a new item is added in one place.

\begin{table}[ht]
\centering\small
\begin{tabularx}{\linewidth}{@{}>{\raggedright\arraybackslash}p{0.17\linewidth} >{\raggedright\arraybackslash}p{0.32\linewidth} X@{}}
\toprule
extension & where it enters & what to provide \\
\midrule
material & \anchor{fea_engine.material.CONSTITUTIVE_REGISTRY}; the \code{D_...} functions & a dataclass and a function returning the constitutive matrix (or a state update for inelastic laws) \\
geometry / mesh & \anchor{fea_engine.geometry.shapes.GEOMETRY_REGISTRY} & a generator returning a \code{Mesh} \\
element & \anchor{fea_engine.elements.ELEMENT_REGISTRY} & an \code{Element} subclass (\Cref{lst:ar-bar}) \\
load & \anchor{fea_engine.loads.LoadPattern}, \code{TimeHistoryLoad}, \code{HarmonicLoad}, \code{PSDLoad} & an object that returns a force vector, for a time or a frequency \\
solution strategy & the drivers of \code{nonlinear_solver}, \code{NewtonOptions} & a function that calls assembly and the linear solve \\
backend & \code{backend=} of \code{FESystem}, \code{backend_dispatch} & \code{"scipy"}, \code{"torch"} or \code{"auto"} (\Cref{sec:ar-backend}) \\
\bottomrule
\end{tabularx}
\caption{Where each kind of extension enters \fea{}.}\label{tab:ar-points}
\end{table}

\section{Does the interface hold? A conformance test}\label{sec:ar-conform}

A registry is only useful if every entry honours the interface. The test of this chapter loops over all \cdeRegTotal{} entries of \code{ELEMENT_REGISTRY} and checks, without needing a material, that each is a subclass of \code{Element}, that \code{n_nodes}, \code{dofs_per_node}, \code{dim} and \code{gauss_order} are positive integers, and, for the interpolated elements, that the shape functions at the natural origin sum to one and that the derivative array has the shape $(\dim,n_{nodes})$ (beam and shell elements, whose first nodal values need not be a partition of unity, skip the sum test). For the contact elements, which have no interpolation, it checks that \code{internal_force} and \code{tangent_stiffness} are callable. \Cref{tab:ar-registry} gives the result: all \cdeRegPass{} entries pass. The test checks the interface and not the physics; the physics is the subject of Chapters~\ref{ch:elem}, \ref{ch:nonlin} and \ref{ch:contact}.

\begin{table}[ht]
\centering\small
\begin{tabular}{lrr}
\toprule
family & registered & pass the interface test \\
\midrule
plane stress / strain & 6 & 6 \\
solid & 11 & 11 \\
plate, shell & 4 & 4 \\
scalar & 1 & 1 \\
beam, truss & 7 & 7 \\
contact & 5 & 5 \\
\midrule
total & 34 & 34 \\
\bottomrule
\end{tabular}

\caption{Interface conformance of the registered elements, by family.}\label{tab:ar-registry}
\end{table}

\section{Two packages, one direction of dependence}\label{sec:ar-packages}

\Cref{tab:ar-packages} gives the size of the packages at the release of this book. The rule that \rom{} works on plain arrays and does not depend on \fea{} can be checked in the source: parsing every module of both packages for import statements shows \cdeFeaImportsRom{} imports of \rom{} in \fea{} and \cdeRomImportsFea{} imports of \fea{} in \rom{}. The tests of \rom{} do import \fea{} (in \cdeRomTestsImportFea{} of \cdeRomTestFiles{} test files), because that is where a reduced model is compared with the full model it came from. The consequence is that \rom{} can reduce a model from any finite element code, as long as the code provides $\bm M$, $\bm K$ and the force as arrays, and that the hooks for hyper-reduction (Chapters~\ref{ch:batch} and~\ref{ch:hyper}) take plain callables instead of an \fea{} object.

\begin{sidenote}{Why the direction of a dependence matters}
A package \emph{depends} on another if it imports it. Then it cannot be used or tested without the other, and a change in the other can break it. \rom{} works on the arrays $\bm M$, $\bm K$ and $\bm F$ and the snapshots of a solution, and never needs to know which program made them, so the same reduction can be applied to the output of \fea{}, of a commercial code, or of an experiment. This is the reason for the one-way rule checked below, and it is also what lets the two packages be installed, versioned and replaced separately.
\end{sidenote}

\begin{table}[ht]
\centering\small
\begin{tabular}{lrrrr}
\toprule
 & source modules & source lines & test files & test lines \\
\midrule
\code{fea_engine} & 62 & 30\,099 & 135 & 27\,703 \\
\code{rom_engine} & 61 & 17\,013 & 85 & 14\,866 \\
\bottomrule
\end{tabular}

\caption{Size of the packages: Python modules and lines (including comments and docstrings) in the source tree and in the tests.}\label{tab:ar-packages}
\end{table}

\section{Backend dispatch and what it costs}\label{sec:ar-backend}

\begin{sidenote}{Backend}
A \emph{backend} is the library that does the dense and sparse linear algebra, NumPy and SciPy by default and PyTorch optionally. The \code{auto} setting chooses between them for each call, and the choice itself has a cost, which this section measures.
\end{sidenote}

\begin{sidenote}{Sparse, dense, direct and iterative}
A \emph{dense} matrix stores all $n^2$ entries; a \emph{sparse} one stores only the nonzeros, which is what makes systems of hundreds of thousands of unknowns possible. A \emph{direct} solver factorises the matrix (Cholesky, $\bm K=\bm L\bm L^T$, if it is symmetric positive definite, LU otherwise) and then solves by substitution; the result is exact up to rounding, and the cost and memory grow quickly with the problem size. An \emph{iterative} solver such as the conjugate-gradient (CG) method only needs products of the matrix with a vector and improves an approximate solution step by step; it is cheap per step, and the number of steps needed grows with the \emph{condition number} of the matrix, the ratio of its largest to its smallest eigenvalue. A matrix is symmetric positive definite (SPD) if $\bm x^T\bm K\bm x>0$ for every nonzero $\bm x$, which is the case for a properly supported elastic structure and is what CG requires.
\end{sidenote}

\code{FESystem(..., backend=...)} chooses the linear-algebra backend of \code{solve_static}: \code{"scipy"} (the default), \code{"torch"} (conjugate gradients on CPU or GPU) or \code{"auto"}. The automatic choice is made at solve time by \anchor{fea_engine.backend_dispatch.select_backend}: a free-dof block that is symmetric positive definite (SPD) and larger than \cdeDispThreshold{} dofs is routed to the torch conjugate-gradient path, everything else to SciPy; a non-SPD matrix is never given to the conjugate-gradient solver. The threshold is a documented rule of thumb and not a benchmark result.

\begin{algbox}{Choice of the linear-algebra backend in \texttt{solve\_static}}
\begin{pseudo}
\Require free-dof block $\bm K_{ff}$, load, \texttt{backend} $\in\{$scipy, torch, auto$\}$
\If{\texttt{backend} = auto}
  \State $\mathrm{spd}\gets$ \texttt{\_is\_spd}$(\bm K_{ff})$ \Comment{dense Cholesky of a densified copy}
  \If{$\mathrm{spd}$ \textbf{and} $n_{ff}>5000$}
    \State \texttt{backend} $\gets$ torch (conjugate gradients)
  \Else
    \State \texttt{backend} $\gets$ scipy (direct)
  \EndIf
\EndIf
\State solve $\bm K_{ff}\bm u_f=\bm f_f$ with the chosen backend
\Ensure displacement; a non-SPD matrix is never sent to conjugate gradients
\end{pseudo}
\end{algbox}

The probe costs $O(n^3)$ time and $O(n^2)$ memory on \emph{every} solve, including those below the threshold, where its answer decides nothing. For a large sparse system the probe and not the solve therefore dominates (measured below), and \code{backend="scipy"} is the better choice there. The code is \anchor{fea_engine.backend_dispatch.select_backend}, and \Cref{fig:ar-dispatch} is the check.

\paragraph{Measured cost of the automatic choice.} \Cref{fig:ar-dispatch} compares \code{solve_static} with \code{backend="scipy"} and \code{backend="auto"} for a sparse plane-stress plate with up to \cdeDispMaxDof{} dofs, a range in which \code{auto} always picks the SciPy path. The two give the same answer, but \code{auto} is slower, and by an amount that grows quickly: at the largest size \code{auto} takes \cdeDispExtraFracLarge{}\% more time than \code{scipy}, \cdeDispExtraMs{} ms more in absolute terms, and the extra time grows with the dof count with an observed exponent of \cdeDispSlope{}. The cause is visible in the source. The SPD test of \code{select_backend} converts the sparse matrix to a dense array and runs a dense Cholesky factorisation, which is $O(n^3)$ in time and $O(n^2)$ in memory, and it does so on every solve, even when the answer decides nothing (below the threshold the SciPy path is taken regardless). For a free-dof count of 5\,000, the dense array takes about \cdeDenseFiveK{} GB; for 50\,000 it would take about \cdeDenseFiveK0{} GB (computed as $8n^2$ bytes, not measured). Until the probe works on the sparse matrix, the practical advice is to pass \code{backend="scipy"} for large sparse systems and to use \code{"auto"} only where the system is small.

\begin{figure}[ht]
\centering
\includegraphics[width=0.95\linewidth]{figures/ch01_dispatch.pdf}
\caption{\code{solve_static} time on a sparse plane-stress plate with \code{backend="scipy"} and \code{backend="auto"} (left), and the extra time of \code{auto} (right). Minimum of three runs on \machine{}. The torch path was not run (torch is \cdeHasTorch{} in the reference setup), and none of the sizes exceeds the threshold at which \code{auto} would choose it.}\label{fig:ar-dispatch}
\end{figure}

\begin{table}[ht]
\centering\small
\begin{tabularx}{\linewidth}{@{}>{\raggedright\arraybackslash}p{0.38\linewidth} X >{\centering\arraybackslash}p{0.11\linewidth}@{}}
\toprule
Concept & Code & Script calls \\
\midrule
system, assembly, boundary conditions, solvers & \anchor{fea_engine.solver.FESystem} & yes \\
element base class and quadrature & \anchor{fea_engine.elements.base.Element}, \anchor{fea_engine.elements.base.gauss_product} & part \\
registries & \anchor{fea_engine.elements.ELEMENT_REGISTRY}, \anchor{fea_engine.geometry.shapes.GEOMETRY_REGISTRY}, \anchor{fea_engine.material.CONSTITUTIVE_REGISTRY} & -- \\
dispatcher & \anchor{fea_engine.backend_dispatch.select_backend} & no \\
\bottomrule
\end{tabularx}
\caption{Where the concepts of this chapter live in the repository (tag \booktag). The last column says whether this chapter's own checks call, by name, the code listed in the row: yes (every name), part, or no. Code reached only through another routine counts as no, so the column is a lower bound on what this book exercises.}\label{tab:ar-map}
\end{table}

\begin{limitsbox}
\begin{itemize}
\item \textbf{The interface test is shallow.} It checks attributes, shape functions and the presence of the force and tangent methods, not that an element is correct (\Cref{tab:ar-registry}).
\item \textbf{Extension is easy for the common case.} A new element that uses the generic Gauss loop is about twenty lines; elements with history variables, mixed formulations or batched kernels have to follow the conventions of Chapters~\ref{ch:nonlin} and~\ref{ch:batch}. Registries are Python dictionaries, so an extension is made in code and there is no plug-in discovery.
\item \textbf{The \code{auto} backend} is costly on sparse systems (\Cref{fig:ar-dispatch}), and its torch branch was not run here.
\item \textbf{Line counts} include comments and docstrings, so they measure documentation as well as code.
\end{itemize}
\end{limitsbox}

\section{Summary}\label{sec:ar-summary}

We started from a request that, in many codes, means editing every solver: a new element. In \fea{} the material, mesh, element, system, solver and post-processing layers are separate, and a solver needs only the matrix, or the force and tangent, from an element. We tested what follows from that. A twenty-line element ran in the existing static and modal solvers with a tip-displacement error of \cdeExtStatic{} without a change to the package source, all \cdeRegTotal{} registered elements pass the interface test, and \rom{} does not depend on \fea{} at source level.

The \code{auto} backend shows why design claims are worth testing: it works, but it is not free, and its cost on sparse systems was not in the documentation. The next chapters use this structure: first the meshes that elements are placed on, and then the elements themselves.

\section*{Exercises}\label{sec:ar-exercises}
\begin{enumerate}
\item \emph{(Implementation)} Add a two-node element for a linear spring ($\bm K=k[[1,-1],[-1,1]]$) as in \Cref{lst:ar-bar}, using \code{stiffness} directly instead of \code{B_matrix}. Which of the interface checks of \Cref{sec:ar-conform} does it still need to pass?
\item \emph{(Implementation)} Extend the conformance test to call \code{stiffness} with a suitable material for every non-contact element and check symmetry and the number of zero eigenvalues (rigid-body modes). Which elements need which material argument?
\item \emph{(Implementation)} Reproduce \Cref{fig:ar-dispatch} with the dense Cholesky of the probe timed on its own. Does it account for the extra time of \code{auto}?
\item \emph{(Implementation)} Write a one-line change of \code{_is_spd} that tests a sparse matrix without densifying it (for example with a sparse Cholesky or an \code{eigsh} test of the smallest eigenvalue), and compare the time with \Cref{fig:ar-dispatch}.
\item \emph{(Interpretation)} Parse the source of \rom{} for the modules that mention \fea{} in comments or docstrings. Is any of them a dependency in disguise (a hidden convention)?
\item \emph{(Derivation)} The probe of the backend dispatcher factors a dense matrix. Show that its cost grows as $n^3/3$ floating-point operations, while a sparse banded Cholesky factorisation costs of the order of $nb^2$ for bandwidth $b$, and find, for $b=50$, the size $n$ above which the probe costs more than the solve it guards.
\end{enumerate}

